\chapter{Conclusion}

In the challenging field of optimization under uncertainty, it is crucial
to guarantee robustness as well as retain objective performance.
\emph{Multi-stage MPC} can handle low-order uncertainty sets but scales
poorly with growing number of scenarios \cite{lucia2015}. This is especially challenging in 
chemical reaction networks, with a high amount of uncertain kinetic parameters
for all appearing reactions. Instead of considering possible parameter combinations online,
neural networks can be trained to capture statistical quantiles offline.
Thus, the computational effort for uncertainty quantification (UQ) is separated from the recursive
optimization. 
\newline
\newline
\noindent
Within the course of this work, several methods have been combined
to derive, implement and assess a NARX framework.
Relying on conformalized quantile regression, an ensemble of neural network NARX
models not only predicts the expected value of the measurement propagation but also
the quantiles of predefined uncertain measurements. Further, a physics-constrained
activation function is presented as a direct corollary (cor.~\ref{cor: physics constrained solution})
of the \emph{Karush-Kuhn-Tucker} conditions. It maps NN outputs during forward passes as well as gradients 
during backward passes into the feasible space defined by linear equality constraints.
The method is demonstrated using a case study
of a 1D-pseudo-homogenious packed-bed tubular reactor model performing the synthesis of ethylene oxide.
Its kinetic parameters are subject to Gaussian distributed uncertainty.
\newline
\newline
\noindent
During NN training, this activation has shown
to significantly speed up covergence, because it incorporates 
equations directly into the network, that govern the data.
PC-NNs with only \SI{25}{\%} of the hidden neurons
of a basic NN exhibit equal validation losses after $512$ epochs.
After excessive training with perfect data, the advantage is assumed to be minor.  
But for small or noisy data sets, PC-NNs portray a real advantage over
basic NNs \cite{chen2024}.
In the forward simulation, physics-consistency is achieved up to machine
precision for both the naive and the the PC-NARX frameworks. The activation
works successfully in training as well as in recursive surrogate model propagation.
\newline
\newline
\noindent
The conditioning of models to quantile predictions is revealed as the Achillis heel of the NARX framework. 
These models exhibit a violation of the conervation quantity of \SI{1.27e-3}{} 
that is twice the value of the expected value model of \SI{6.39e-4}{}. 
This is anticipated, because the distance-based weighting (def.~\ref{sur:def: gaussian weight function})
focuses on data points, that are 
closest to the quantile prediction. Consequently, the considered data is reduced dramatically.
Further, the conditioned quantile predictions cross the expected value trajectory. Neither CQR nor
a PC activation function can provide a remedy. This underlines, that the conditioning of 
expected value models with according quantile predictions is only a pragmatic implementation
to allow a recursive simulation. While it captures some correlation
between measurements, it lacks rigorousitiy. Quantile predictions are only valid when calculated
in a single-step. The naive forward simulation represents an approximation of the quantiles at best.
This is prooven by the fact, that the desired coverage of $0.8$ is only achieved at the first
measurement position. Neither the vanilla, naive nor the physics-constrained ensemble
yield a significant change in mean coverage nor relative interval width. Solely, the PC framework
achieves more reproducable coverages with a standard deviation that is \SI{12.9}{\%} lower 
than the vanilla framework on average. All versions show increasing over-coverage 
over the reactor length. The heteroscedasticity of the plug-flow behavior
with uncertain kinetics is not learned apropriately.
\newline
\newline
\noindent
To counteract these downsides, two improvements should be made. First, multi-step NARX models
ensure correct quantiles for the complete prediction horizon. This replaces the recursive
single-step simulation at the cost of a greater network size. Second, the measurement position
should act as an input to the NN to learn a complete reactor profile and reduce the NN output size.
During inference, the
measurement position can be held constant. Additionally, a greater amount of data is available, as 
all spatial discretization nodes can be considered during training.
In the mere forward pass, the physics-constrained
activation function is another useful tool for surrogate modelling with cosmetic character, 
whereas it cannot solve errors introduced by foundational assumptions in the NARX architecture.
\newline
\newline
\noindent
For the application within an MPC algorithm, performance, constraint violation, control effort and 
runtime are evaluated.
The uncertainty-aware MPCs enable safe operation with zero temperature limit violation
which is impressive, as the heat conductivity of the packed-bed can easily vary between
\SI{0.3}{} and \SI{0.5}{\watt \; \meter^{-1} \; \kelvin^{-1}} (fig.~\ref{case:def: heat conductivity distribution}).
While the performance stays unaffected in this specific example,
UQ based frameworks reduce the combined constraint violation by \SI{95.0}{\%} compared to the nominal case. 
Because of the increased complexity, the solution time grows by \SI{216}{\%} here.
Furthermore, input oscillations shrink by \SI{39.2}{\%} due to a reasonable safety-distance
to the critical contraint.
CQR is a useful tool to ensure robustness up to a specified coverage and enables 
to capture high-order uncertainties. However, the number of uncertain measurements to be accounted
for remains a limitation, as one separate quantile regressor needs to be trained for each.
\newline
\newline
\noindent
The integration of a PC activation function enables an improvement in the
conversion constraint violation of \SI{16.2}{\%} relative to the other UQ architectures
and control effort of \SI{44.4}{\%} relative to the nominal one.
This is achieved neither by sacrificing objective performance, nor exploiting infeasible
temperatures. Additionally, the solution time only increases by \SI{151.7}{\%} to \SI{3.80}{\second}.
The reduction of input oscillations, faster convergence in MPC operation and in NN training
is not only caused by correcting the NN outputs themselves. The adjustment of the gradient
during the backward passes (eq.~\ref{results:eq:pc-effect-on-backpropagation}) induce the major impact.
This type of activation function allows to include linear equality constraints into machine learning
architectures and is not only useful for chemical engineering.
It is a valuable modeling tool with a wide range of possible applications. 